\section*{Chapter \ref{ch:b1:organic-redox} --- Oxidation and Reduction in Organic Chemistry}

\begin{solution}{exo:b1:organic-redox:1}
Ethanol: \ce{CH3} $-3$, \ce{CH2OH} $-1$. Ethanal: \ce{CH3} $-3$, CHO $+1$.
Ethanoic acid: \ce{CH3} $-3$, COOH $+3$. Propanone: \ce{CH3} $-3$ (twice),
C=O $+2$. Methanoic acid: $+2$.
\end{solution}

\begin{solution}{exo:b1:organic-redox:2}
Propan-1-ol: propanal, then propanoic acid. Propan-2-ol: propanone.
2-Methylpropan-2-ol: no reaction (the orange colour stays).
\end{solution}

\begin{solution}{exo:b1:organic-redox:3}
\ce{CH3COCH3 + 2H+ + 2e- -> CH3CH(OH)CH3};
\ce{CH3CHO + 2H+ + 2e- -> CH3CH2OH}.
\end{solution}

\begin{solution}{exo:b1:organic-redox:4}
\ce{NaBH4}: butan-1-ol, butan-2-ol, no reaction, no reaction (the acid is
only deprotonated). \ce{LiAlH4}: butan-1-ol, butan-2-ol, butan-1-ol and
ethanol, butan-1-ol.
\end{solution}

\begin{solution}{exo:b1:organic-redox:5}
\ce{5CH3CH(OH)CH3 + 2MnO4- + 6H+ -> 5CH3COCH3 + 2Mn^2+ + 8H2O}.
$n = 1.20/60.0 = \qty{0.0200}{mol}$; permanganate $\frac25 \times 0.0200 =
\qty{8.0e-3}{mol}$, \qty{400}{mL} of the \qty{0.020}{mol/L} solution.
\end{solution}

\begin{solution}{exo:b1:organic-redox:6}
Heat the alcohol with the \omterm{def:b1:nernst:couple}{oxidant} in a flask fitted for distillation, at a
temperature between the two boiling points: propanal (\qty{48}{\celsius})
distils as it forms and escapes further oxidation; propan-1-ol
(\qty{97}{\celsius}) stays in the flask.
\end{solution}

\begin{solution}{exo:b1:organic-redox:7}
A B--H pair of \ce{BH4-} attacks the carbonyl carbon, the C=O $\pi$ pair
going to oxygen; the \omterm{def:b1:alcohol-activation:alkoxide}{alkoxide} takes a proton from ethanol. The hydrogen
on the carbon of the product (\ce{CH3CH2OH}, one of the two on C1) comes
from the reagent; the one on oxygen from the solvent.
\end{solution}

\begin{solution}{exo:b1:organic-redox:8}
Butan-2-ol is \omterm{def:b1:stereochemistry-in-depth:stereocentre}{chiral}, but the planar ketone is attacked equally on both
faces: a \omterm{def:b1:stereochemistry-in-depth:enantiomers}{racemic mixture}, optically inactive.
\end{solution}

\begin{solution}{exo:b1:organic-redox:9}
\ce{3CH3CH2OH + 2Cr2O7^2- + 16H+ -> 3CH3COOH + 4Cr^3+ + 11H2O}: orange
dichromate (chromium $+$VI) becomes green chromium(III).
\end{solution}

\begin{solution}{exo:b1:organic-redox:10}
2-Methylbutan-2-ol: ethanal + \ce{CH3CH2MgBr} gives butan-2-ol (an
addition, no redox); oxidation to butanone; \ce{CH3MgBr} on butanone,
then work-up. One oxidation. Pentan-3-ol: propanal + \ce{CH3CH2MgBr} in
one step, no redox.
\end{solution}

\begin{solution}{exo:b1:organic-redox:11}
1. Protect the ketone as a cyclic \omterm{def:b1:acetals-protection:hemiacetal}{acetal} (ethane-1,2-diol, acid,
Dean--Stark). 2. \ce{LiAlH4} in dry ether reduces the ester to the
primary alcohol (and methanol). 3. Aqueous acid removes the \omterm{def:b1:acetals-protection:hemiacetal}{acetal}:
5-hydroxypentan-2-one, \ce{CH3CO(CH2)3OH}.
\end{solution}

\begin{solution}{exo:b1:organic-redox:12}
\ce{RCH2OH}: carbinol carbon at $-1$; \ce{RCOOH}: $+3$; four electrons, as in
\ce{RCH2OH + H2O -> RCOOH + 4H+ + 4e-}. Methane ($-4$) to carbon dioxide
($+4$): eight electrons, \ce{CH4 + 2H2O -> CO2 + 8H+ + 8e-}.
\end{solution}

\begin{solution}{pb:b1:organic-redox:1}
\textbf{1.} C1 (bearing OH): 0; the five \ce{CH2}: $-2$.
\textbf{2.} C1 (C=O): $+2$; the five \ce{CH2}: $-2$.
\textbf{3.} The two COOH carbons: $+3$; the four \ce{CH2}: $-2$.
\textbf{4.} C1 ($0 \to +3$) and its neighbour C6 ($-2 \to +3$), which
becomes the second carboxylic carbon; the C1--C6 bond is broken.
\textbf{5.} $3 + 5 = 8$ electrons.
\textbf{6.} Nitric acid is strong enough to break the C--C bond next to the
carbonyl, which milder \omterm{def:b1:nernst:couple}{oxidants} do not do.
\textbf{7.} \ce{C6H12O + 3H2O -> C6H10O4 + 8H+ + 8e-}.
\textbf{8.} $+$V and $+$I.
\textbf{9.} \ce{2HNO3 + 8H+ + 8e- -> N2O + 5H2O}.
\textbf{10.} Adding, eight electrons, \ce{8H+} and three water molecules
cancel: \ce{C6H12O + 2HNO3 -> C6H10O4 + N2O + 2H2O}.
\textbf{11.} 8.
\textbf{12.} $10^6/146.0 = \qty{6.85e3}{mol}$ of acid; $2 \times 6.85 \times
10^3 \times 63.0 = \qty{8.6e5}{g}$, \qty{0.86}{t} of nitric acid.
\textbf{13.} \ce{H2O2 + 2H+ + 2e- -> 2H2O}.
\textbf{14.} \ce{C6H12O + 4H2O2 -> C6H10O4 + 5H2O}.
\textbf{15.} $4 \times 6.85 \times 10^3 \times 34.0 = \qty{9.3e5}{g}$,
\qty{0.93}{t}.
\textbf{16.} Water: no nitrogen oxide.
\textbf{17.} A large equilibrium constant says nothing about the rate;
\ce{H2O2} reacts slowly without a \omterm{def:b1:elementary-steps:catalyst}{catalyst}, as its own decomposition does.
\textbf{18.} It reduces cyclohexanone back to cyclohexanol:
\ce{4C6H10O + BH4- -> B(OC6H11)4-}, then
\ce{B(OC6H11)4- + 4H2O -> 4C6H11OH + B(OH)4-}.
\textbf{19.} \qty{6.85e3}{mol}.
\textbf{20.} \qty{6.85e3}{mol} of \ce{N2O} (one per acid).
\textbf{21.} $6.85 \times 10^3/0.95 \times 100.0 = \qty{7.2e5}{g}$, \qty{0.72}{t}.
\textbf{22.} $6.85 \times 10^3 \times 44.0 = \qty{3.0e5}{g}$: \textbf{\qty{0.30}{t}
of \ce{N2O} per tonne of adipic acid}.
\end{solution}
